\documentclass[aspectratio=169,10pt]{beamer}
\newcommand{\TalkShort}{Paper review: Neggers et al.\ (2003)}
% Visual reference: Matteo Di Sante, 13_06_26_slides_soaring_ctrw.tex.
% Steel-blue title bar, Palatino, Pisa seal, white pages and a pale title band.
\usepackage[T1]{fontenc}
\usepackage[utf8]{inputenc}
\usepackage[english]{babel}
\usepackage{mathpazo}
\usefonttheme{serif}
\usepackage{amsmath,amssymb,bm,booktabs,array,graphicx,tikz,microtype}
\usetikzlibrary{arrows.meta,calc,positioning}
\usepackage{siunitx}
\sisetup{group-separator={,},group-minimum-digits=4}
\usepackage{pgfpages}
\ifdefined\Speakernotes
  \setbeameroption{show notes on second screen=right}
\else
  \setbeameroption{hide notes}
\fi
\definecolor{pisablue}{RGB}{74,96,121}
\definecolor{pisablight}{RGB}{192,206,220}
\definecolor{accent}{RGB}{26,118,179}
\definecolor{para}{HTML}{3477A8}
\definecolor{hang}{HTML}{B5482A}
\definecolor{beginner}{HTML}{6A3D9A}
\definecolor{expert}{HTML}{C98A1E}
\definecolor{teal}{HTML}{2E7D8A}
\definecolor{wine}{HTML}{8E3B5C}
\definecolor{muted}{HTML}{666666}
\setbeamercolor{structure}{fg=pisablue}
\setbeamercolor{normal text}{fg=black,bg=white}
\setbeamercolor{frametitle}{fg=white,bg=pisablue}
\setbeamercolor{alerted text}{fg=accent}
\setbeamercolor{block title}{fg=white,bg=pisablue}
\setbeamercolor{block body}{fg=black,bg=pisablight!28}
\setbeamercolor{discussion}{fg=pisablue,bg=pisablight!30}
\setbeamercolor{button}{fg=white,bg=pisablue}
\setbeamerfont{button}{size=\normalsize}
\setbeamersize{text margin left=7mm,text margin right=7mm}
\setbeamertemplate{navigation symbols}{}
\setbeamertemplate{itemize item}{\textcolor{pisablue}{\small$\bullet$}}
\setbeamertemplate{itemize subitem}{\textcolor{pisablue}{\textendash}}
\setbeamertemplate{itemize/enumerate body begin}{\setlength{\itemsep}{0.55em}}
\setbeamertemplate{frametitle}{%
  \nointerlineskip
  \begin{beamercolorbox}[wd=\paperwidth,ht=9mm,dp=3mm,center]{frametitle}
    \large\insertframetitle\strut
  \end{beamercolorbox}
  \begin{tikzpicture}[remember picture,overlay]
    \node[anchor=north east,text=white,inner sep=0pt,xshift=-2.3mm,yshift=-4.2mm]
      at(current page.north east){\small\insertframenumber};
  \end{tikzpicture}
}
\setbeamertemplate{footline}{%
  \hspace{7mm}\textcolor{muted}{\fontsize{6.5}{8}\selectfont
    Matteo Di Sante\enspace|\enspace\TalkShort}\vspace{2mm}
}
\newcommand{\kw}[1]{\textcolor{accent}{#1}}
\newcommand{\source}[1]{\par\vspace{1mm}{\fontsize{7}{8.3}\selectfont\color{muted}#1\par}}
% Discussion prompts belong only on the companion notes page.
\newcommand{\question}[1]{\note{\par\textbf{Discussion.} #1\par\medskip}}
\newcommand{\takeaway}[1]{\par\vfill
  \begin{beamercolorbox}[wd=\linewidth,sep=5pt]{discussion}
    \small #1
  \end{beamercolorbox}}
\newcommand{\fig}[2][54mm]{\includegraphics[width=\linewidth,height=#1,keepaspectratio]{assets/#2}\par}
\newcommand{\speaker}[1]{\note{#1}}
\setbeamertemplate{note page}{%
  \vspace{4mm}\noindent%
  \begin{minipage}{\linewidth}
    {\small\color{pisablue}\TalkShort\quad|\quad Slide \insertframenumber}\par\medskip
    {\normalsize\bfseries\insertframetitle}\par\medskip
    {\small\insertnote}
  \end{minipage}
}
\newcommand{\refitem}[3]{\href{#1}{#2}\hfill {\color{muted}#3}\par\medskip}
\newcommand{\titleframe}[4]{%
\begin{frame}[plain]
\begin{tikzpicture}[remember picture,overlay]
 \fill[pisablue](current page.south west)rectangle(current page.north east);
 \fill[pisablight](current page.south west)rectangle([yshift=20mm]current page.south east);
 \node[text=white,align=center,text width=.92\paperwidth]at(current page.center)
   {{\LARGE\bfseries #1}};
 \node[text=pisablue,align=center]at([yshift=10mm]current page.south)
   {{\normalsize Version: 14/09/2026}\\[2mm]
    {\small #3}};
\end{tikzpicture}
\note{#4}
\end{frame}}
% Marks a thesis subsection boundary: a PDF bookmark plus a full-page divider,
% so each block of slides can be traced back to its exact place in the thesis.
\newcommand{\subsectionframe}[2]{%
\subsection{#2}
\begin{frame}[plain,noframenumbering]
\begin{tikzpicture}[remember picture,overlay]
 \fill[pisablue](current page.south west)rectangle(current page.north east);
 \node[text=pisablight,align=center,text width=.82\paperwidth]
   at([yshift=10mm]current page.center){{\normalsize Thesis~\S#1}};
 \node[text=white,align=center,text width=.82\paperwidth]
   at([yshift=-4mm]current page.center){{\Large\bfseries #2}};
\end{tikzpicture}
\end{frame}}
\newcommand{\backupstart}{%
\appendix
\begin{frame}[plain,noframenumbering]
\begin{tikzpicture}[remember picture,overlay]
 \fill[pisablue](current page.south west)rectangle(current page.north east);
 \node[text=white,align=center]at(current page.center)
  {{\LARGE\bfseries Discussion material}\\[4mm]{\large Extra figures, derivations and references}};
\end{tikzpicture}
\end{frame}}
\hypersetup{colorlinks=true,urlcolor=accent,linkcolor=accent,pdfauthor={Matteo Di Sante}}
% A notes sheet contains a second rendering of the slide. Disable its duplicate
% destinations; the normal slide PDF retains all hyperlinks and the viewer button.
\ifdefined\Speakernotes\hypersetup{draft}\fi

% Shared layout for the paper-review decks. It reuses ../theme.tex and the
% absolute-position slide macros of offsite2026/offsite-2026.tex (160 x 90 mm).
\setbeamertemplate{frametitle}{}
\setbeamertemplate{footline}{}
\setbeamerfont{note page}{size=\small}
\newcommand{\txt}[5][]{\node[anchor=north west,inner sep=0pt,text width=#4mm,align=left,font=\fontsize{10}{12}\selectfont,#1] at (#2,#3) {#5};}
\newcommand{\smalltxt}[5][]{\txt[font=\fontsize{8}{10}\selectfont,#1]{#2}{#3}{#4}{#5}}
\newcommand{\tinytxt}[5][]{\txt[font=\fontsize{7}{8.6}\selectfont,#1]{#2}{#3}{#4}{#5}}
\newcommand{\head}[4]{\txt[text=pisablue,font=\bfseries\fontsize{10}{12}\selectfont]{#1}{#2}{#3}{#4}}
% \figat{x}{y}{max width}{max height}{file}: top-left anchored, aspect kept.
\newcommand{\figat}[5]{\node[anchor=north west,inner sep=0pt]at(#1,#2){\includegraphics[width=#3mm,height=#4mm,keepaspectratio]{#5}};}
\newcommand{\bottomline}[1]{\fill[pisablight!33](7,75.5)rectangle(153,83);\txt[text=pisablue,font=\fontsize{10}{12}\selectfont]{10}{77}{140}{#1}}
\newcommand{\credit}[1]{\txt[text=muted,font=\fontsize{6}{7}\selectfont]{7}{86}{146}{#1}}
\newcommand{\startslide}[1]{\begin{tikzpicture}[remember picture,overlay,shift={(current page.north west)},x=1mm,y=-1mm]
 \fill[white](0,0)rectangle(160,90);
 \fill[pisablue](0,0)rectangle(160,12.5);
 \txt[text=white,font=\fontsize{15}{17}\selectfont\bfseries]{7}{3.2}{143}{#1}
 \node[anchor=east,inner sep=0pt,text=white,font=\fontsize{8}{9}\selectfont]at(155,6.6){\insertframenumber};
}
\newcommand{\finishslide}{\end{tikzpicture}}
% \lead{text}: one or two plain sentences under the title, saying what the slide shows.
\newcommand{\lead}[1]{\txt[font=\fontsize{9.5}{11.6}\selectfont]{7}{15}{146}{#1}}
% \symbox[title]{x}{y}{width}{\sym{..}{..}...}: legend of every symbol used on the slide.
\newcommand{\symbox}[5][Symbols]{\node[anchor=north west,inner sep=1.6mm,fill=pisablight!18,draw=pisablue!45,line width=.3pt,text width=#4mm,align=left,font=\fontsize{7.8}{9.2}\selectfont]at(#2,#3){{\color{pisablue}\bfseries #1}\par\vspace{.2mm}#5};}
% \symboxtwo{x}{y}{left width}{right width}{left \sym list}{right \sym list}: two-column legend.
\newcommand{\symboxtwo}[6]{\node[anchor=north west,inner sep=1.6mm,fill=pisablight!18,draw=pisablue!45,line width=.3pt,align=left,font=\fontsize{7.8}{9.2}\selectfont]at(#1,#2){{\color{pisablue}\bfseries Symbols}\\[-2.6mm]\begin{tabular}{@{}p{#3mm}@{\hspace{4mm}}p{#4mm}@{}}#5&#6\end{tabular}};}
\newcommand{\sym}[2]{\par\vspace{.5mm}\hangindent=3mm\hangafter=1\noindent#1\enspace{\color{black!75}#2}}
% \figcap{x}{y}{width}{text}: muted caption under a figure.
\newcommand{\figcap}[4]{\txt[text=muted,font=\fontsize{7.2}{8.6}\selectfont]{#1}{#2}{#3}{#4}}
% Two kinds of boxed statement: what the paper says, and what we propose.
% The node width is #3 mm plus 3 mm of padding.
\newcommand{\paperbox}[4]{\node[anchor=north west,inner sep=1.5mm,fill=pisablight!22,draw=pisablue!55,line width=.4pt,text width=#3mm,align=left,font=\fontsize{7.5}{9.3}\selectfont]at(#1,#2){{\color{pisablue}\bfseries In the paper}\\[.8mm]#4};}
\newcommand{\ideabox}[4]{\node[anchor=north west,inner sep=1.5mm,fill=expert!9,draw=expert,line width=.5pt,text width=#3mm,align=left,font=\fontsize{7.5}{9.3}\selectfont]at(#1,#2){{\color{expert!70!black}\bfseries Proposal, not in the paper}\\[.8mm]#4};}
\newcommand{\ourtag}{{\color{expert!70!black}\bfseries[ours]}}
% \coverslide{title}{authors}{journal}{question}{tikz decoration}{speaker note}
\newcommand{\coverslide}[6]{%
\begin{frame}[plain]
\begin{tikzpicture}[remember picture,overlay,shift={(current page.north west)},x=1mm,y=-1mm]
 \fill[pisablue](0,0)rectangle(160,90);
 \fill[pisablight](0,73)rectangle(160,90);
 \txt[text=white,font=\fontsize{21}{25}\selectfont\bfseries]{11}{12}{140}{#1}
 \txt[text=pisablight,font=\fontsize{12}{15}\selectfont]{11}{37}{138}{#2}
 \txt[text=white,font=\fontsize{9.5}{12}\selectfont]{11}{44}{138}{#3}
 \txt[text=pisablight,font=\itshape\fontsize{10}{12.5}\selectfont]{11}{52}{80}{#4}
 #5
 \txt[text=pisablue,font=\fontsize{12}{14}\selectfont]{11}{78}{70}{Matteo Di Sante}
 \txt[text=pisablue,font=\fontsize{9}{11}\selectfont,align=right]{88}{78.5}{61}{Paper review}
\end{tikzpicture}
\note{#6}
\end{frame}}
\newcommand{\notetext}[1]{\note{#1}}
% Denser note pages than the offsite deck: the review notes explain whole sections.
\setbeamertemplate{note page}{%
  \vspace{3mm}\noindent%
  \begin{minipage}{\linewidth}
    {\fontsize{7}{8.4}\selectfont\color{pisablue}\TalkShort\quad|\quad Slide \insertframenumber}\par\smallskip
    {\small\bfseries\insertframetitle}\par\smallskip
    {\fontsize{7.3}{8.8}\selectfont\insertnote}
  \end{minipage}
}

% Diagram of the spatial observables sought in this review. Illustrative only.
\newcommand{\thermalmap}{%
 \smalltxt[text=muted]{7}{25}{70}{Spatial quantities (schematic)}
 \draw[-{Latex},muted,line width=.65pt](12,68)--(77,68);
 \draw[-{Latex},muted,line width=.65pt](12,68)--(12,34);
 \node[anchor=west,text=muted,font=\fontsize{10}{12}\selectfont]at(77,68){$x$};
 \node[anchor=east,text=muted,font=\fontsize{10}{12}\selectfont]at(11,34){$y$};
 \draw[pisablue,line width=.9pt,fill=pisablight!20](30,46)circle(7);
 \draw[pisablue,line width=.9pt,fill=pisablight!20](63,46)circle(10);
 \draw[<->,wine,line width=.8pt](30,46)--(63,46);
 \node[anchor=south,text=wine,font=\fontsize{11}{13}\selectfont]at(46.5,44){$d_{12}$};
 \draw[-{Latex},pisablue,line width=.8pt](30,46)--(30,39);
 \node[anchor=south,text=pisablue,font=\fontsize{10}{12}\selectfont]at(30,38){$R_1$};
 \fill[pisablue](30,46)circle(.7);
 \fill[pisablue](63,46)circle(.7);
 \txt[align=center,text=pisablue]{19}{56}{23}{$(x_1,y_1)$}
 \txt[align=center,text=pisablue]{52}{59}{24}{$(x_2,y_2)$}
}

\hypersetup{pdftitle={Neggers et al. (2003): cloud-size statistics},pdfsubject={Five-slide paper review with one appendix slide}}
\newcommand{\figdir}{assets/neggers2003}
\begin{document}

% 1: measured objects and definitions ---------------------------------
\begin{frame}[plain]\frametitle{Neggers et al. (2003): cloud-size statistics}\startslide{Neggers et al.\ (2003): cloud-size statistics}
 \lead{Large-eddy simulations (LES) of shallow cumulus: simulated clouds, not measurements. Each cloud gets one size, and the paper studies, among other things, the histogram of these sizes.}
 \figat{8}{25}{57}{47}{\figdir/fig3.png}
 \txt[font=\fontsize{15}{18}\selectfont]{69}{26}{40}{$l=\sqrt{A^p}$}
 \txt[font=\fontsize{13}{16}\selectfont]{69}{37}{40}{$N=\displaystyle\int_0^\infty\mathcal N(l)\,dl$}
 \smalltxt{69}{49}{40}{\textbf{Three simulated cases}, set up like real campaigns:
 \begin{itemize}\setlength{\itemsep}{0pt}\setlength{\topsep}{0pt}\setlength{\parsep}{0pt}
 \item BOMEX: Barbados, ocean.
 \item SCMS: Florida, land.
 \item ARM: Oklahoma, land, diurnal cycle.
 \end{itemize}}
 \symbox{110}{25}{40}{%
  \sym{cloud:}{group of touching saturated grid cells in 3D}
  \sym{$A^p$}{vertically projected area of one cloud: its footprint seen from above}
  \sym{$l$}{cloud size, m}
  \sym{$\mathcal N(l)$}{number of clouds per unit size, m$^{-1}$}
  \sym{$N$}{total number of clouds in the domain}}
 \bottomline{Find each cloud in 3D $\;\to\;$ project it vertically $\;\to\;$ histogram of the sizes $l$}
 \credit{Neggers, Jonker \& Siebesma, J. Atmos. Sci. 60, 1060--1074 (2003). \S2--3, Fig. 3, Eqs. (1)--(2). Grid 50 m, domain 6.4 km.}
\notetext{Size Statistics of Cumulus Cloud Populations in Large-Eddy Simulations, Neggers, Jonker and Siebesma (2003). The study measures the sizes of shallow cumulus clouds in LES and their contributions to cloud cover and vertical transport. BOMEX represents steady trade-wind cumulus over the ocean. SCMS and ARM represent continental shallow cumulus, with a diurnal cycle in ARM. All clouds are simulated: each case is initialised and forced with measurements from a real campaign (Barbados, Cocoa Beach in Florida on 5 August 1995, Oklahoma SGP on 21 June 1997), but the simulated clouds are not the observed ones. Comparison with nature is through published exponents and dominant sizes. The reference horizontal grid spacing is 50 m, in a 6.4 km domain. Repeated snapshots sample tens of thousands of clouds per case.\par\smallskip A cloud is a connected set of saturated grid cells. Figure 3 shows the distinction between a cross-section at one height and the vertical projection of the whole cloud. The chosen size is $l=\sqrt{A^p}$, the side of a square with the cloud's projected area, not a radius or a diameter. It is chosen so that LES can be compared with satellite images, which see only projections. The size density $\mathcal N(l)$ gives counts per unit size, and its integral is the total count. The bottom line summarizes the actual measurement procedure. Source: Sections 2--3 and Appendix, Figure 3, Equations (1)--(2).}
\finishslide\end{frame}

% 2: fitted size density -----------------------------------------------
\begin{frame}[plain]\frametitle{A power law below a variable size break}\startslide{A power law below a variable size break}
 \lead{The histogram counts clouds per size: \textbf{small clouds are by far the most numerous}. Below a size $\ell_d$ their number falls as a power law of size; above $\ell_d$ it falls faster.}
 \figat{8}{25}{46}{42}{\figdir/fig4a.png}
 \figat{54.5}{25}{44}{42}{\figdir/fig6.png}
 \figcap{8}{68.5}{46}{Fig. 4a: size density against $l$.}
 \figcap{54.5}{68.5}{44}{Fig. 6: against $l/\ell_d$, collapsed.}
 \txt[font=\fontsize{15}{18}\selectfont]{101}{25}{52}{$\mathcal N(l)=a\,l^{\,b}$}
 \txt{101}{34}{52}{$b=-1.70$\ \ for\ \ $l<\ell_d$}
 \symbox{101}{42}{49}{%
  \sym{$\ell_d$}{scale break: size where the exponent changes. BOMEX 700 m, SCMS 1050 m, ARM 400 $\to$ 1250 m in the afternoon}
  \sym{$a,\,b$}{fit constants; $b$ is the exponent}
  \sym{$\mathcal N^*$}{$=l\ln(10)\,\mathcal N$, clouds per decade of size (plotted), slope $b+1=-0.70$}}
 \bottomline{Same exponent in all cases: only the break $\ell_d$ changes}
 \credit{\S4a, Figs. 4a and 6, Table 2, Eqs. (8)--(9). Plotted: $\mathcal N^*/N$ in log bins.}
\notetext{Figure 4a shows the normalized size density in logarithmic bins. Below the scale break, the three cases have approximately the same slope. The fitted log-bin slope is $-0.70$, corresponding to the exponent $b=-1.70$ of $\mathcal N(l)=a\,l^b$ because $\mathcal N^*=l\ln(10)\mathcal N$: a power law per decade has its exponent raised by one. The supported range spans roughly one decade between the grid scale and the break.\par\smallskip Above $\ell_d$, the density decreases faster. Table 2 gives 700 m for BOMEX and 1050 m for SCMS. During the ARM afternoon, the break increases from 400 to 1250 m. Increasing the number of sampled clouds does not remove the break.\par\smallskip Figure 6 plots the three size densities against $l/\ell_d$. Their approximate collapse supports a common functional form for these simulated cumulus populations, with a case-dependent break: the break is the only variable. Its physical origin remains unresolved; the paper mentions subcloud-layer depth and humidity fluctuations as candidates. It is a cloud-size scale, not an inter-cloud distance. Source: Section 4a, Figures 4a and 6, Table 2, Equations (8)--(9).}
\finishslide\end{frame}

% 3: area weighting ----------------------------------------------------
\begin{frame}[plain]\frametitle{The size break sets the dominant cloud cover}\startslide{The size break sets the dominant cloud cover}
 \lead{Weighting every cloud by its area gives the cloud cover contributed by each size. It grows up to $\ell_d$ and drops beyond, so \textbf{mid-size clouds cover most of the sky.}}
 \txt[font=\fontsize{13}{16}\selectfont]{8}{26}{85}{$\alpha^p(l)=\dfrac{l^2\,\mathcal N(l)}{L_xL_y}\ \propto\ l^{\,b+2}=l^{\,0.30}\quad(l<\ell_d)$}
 \txt[font=\fontsize{13}{16}\selectfont]{8}{39}{85}{$a^p=\displaystyle\int_0^\infty\alpha^p(l)\,dl$}
 \symbox{8}{51}{86}{%
  \sym{$\alpha^p(l)$}{cloud cover from clouds of size $l$, per unit size, m$^{-1}$}
  \sym{$a^p$}{projected cloud cover: fraction of the domain covered, seen from above}
  \sym{$L_xL_y$}{horizontal area of the domain}}
 \figat{102}{24.5}{50}{46.5}{\figdir/fig7a.png}
 \figcap{102}{71.2}{50}{Fig. 7a: $\alpha^p(l)$, three cases.}
 \bottomline{Cloud cover peaks near the break $\ell_d$}
 \credit{\S3--4b, Fig. 7a, Eqs. (3)--(4), (10).}
\notetext{Figure 7a gives the contribution to projected cloud cover per unit size. Each cloud contributes its projected area $l^2$, so $\alpha^p(l)=l^2\mathcal N(l)/(L_xL_y)$, where $L_xL_y$ is the horizontal domain area. For a power law with exponent $b$, the cover decomposition scales as $l^{b+2}$. Below the break, $b=-1.70$ gives a positive exponent, $0.30$: the many small clouds are too small to dominate. Above the break the density declines faster than $l^{-2}$, so the cover decomposition decreases. The peak therefore lies at an intermediate size near the break. Without a break, an exponent above $-2$ would let the largest clouds dominate.\par\smallskip This is a decomposition per unit linear size, not per logarithmic bin. The total projected cover is $a^p=\int_0^\infty\alpha^p(l)\,dl$. The paper notes that real shallow cumulus fields also show this intermediate dominant size. Source: Section 3, Section 4b, Figure 7a, Equations (3)--(4) and (10).}
\finishslide\end{frame}

% 4: vertical transport ------------------------------------------------
\begin{frame}[plain]\frametitle{Mass flux also peaks at intermediate sizes}\startslide{Mass flux also peaks at intermediate sizes}
 \lead{Larger clouds rise faster, so mass flux still peaks at \textbf{intermediate} sizes, near the cover peak, but the smallest clouds drop out: their vertical velocity is very low.}
 \figat{8}{25}{44}{43}{\figdir/fig7b.png}
 \figat{53.5}{25}{43}{43}{\figdir/fig7c.png}
 \figcap{8}{68.5}{44}{Fig. 7b: mean vertical velocity $w$.}
 \figcap{53.5}{68.5}{44}{Fig. 7c: mass flux $\mu(l)$.}
 \txt[font=\fontsize{12.5}{15}\selectfont]{100}{23.5}{53}{{\small Mass flux:}\\$\mu(l,z)=\alpha(l,z)\,w(l,z)$}
 \txt[font=\fontsize{12.5}{15}\selectfont]{100}{34}{53}{$\mu(l)=\dfrac1{h_c}\displaystyle\int_{h_c}\mu(l,z)\,dz$}
 \symbox{99}{44.5}{51}{%
  \sym{$w(l,z)$}{mean vertical velocity of clouds of size $l$ at height $z$, m\,s$^{-1}$}
  \sym{$\alpha(l,z)$}{their cloud cover at height $z$, per unit size}
  \sym{$h_c$}{depth of the cloud layer}
  \sym{$\mu$}{mass flux divided by air density, per unit size, s$^{-1}$}}
 \bottomline{Same intermediate dominant size as cover; the smallest clouds carry almost nothing}
 \credit{\S3--4b, Fig. 7b--c, Eqs. (5)--(7). $l$ is always the projected size, also at a single height $z$.}
\notetext{Figure 7b shows that cloud-mean vertical velocity generally increases with cloud size. Figure 7c shows the resulting mass-flux decomposition. At height $z$, $\mu(l,z)=\alpha(l,z)w(l,z)$, where $\alpha(l,z)$ is the cloud-fraction contribution per unit size at that height. Averaging over the cloud-layer depth $h_c$ gives $\mu(l)$. Its integral over size is the total kinematic mass flux, in m/s under this paper's convention: mass flux divided by air density.\par\smallskip Cloud size is always the projected size, even when evaluating the flux at a specific height, so each cloud stays in one size bin at all heights. Since $w$ increases with size, the small-cloud side is suppressed relative to the cover: the paper describes the dominant size as better defined and shifted somewhat toward larger sizes. Comparing Figs. 7a and 7c, the peak positions are nearly the same (read off the plots, about 650, 1000 and 1500 m for BOMEX, SCMS and ARM), so the shift is in the weight of the distribution, not in the position of its maximum. The smallest clouds contribute very little to total vertical transport, mainly because of their very low vertical velocities.\par\smallskip The intermediate dominant size persists for tested grid spacings of 25--100 m. A 3.2 km domain truncates the largest clouds. Increased shear broadens projected sizes while leaving total mass flux almost unchanged. These sensitivity tests support the main transport result but do not identify the mechanism of the size break. Source: Section 3, Section 4b--e, Figure 7b--c and Figures 9--11.}
\finishslide\end{frame}

% 5: climbable clouds versus thermal diameters ---------------------------
\begin{frame}[plain]\frametitle{``Climbable clouds'' are far wider than thermals}\startslide{``Climbable clouds'' are far wider than thermals}
 \lead{The paper gives no thermal statistics. Fig.\ 7b gives the mean updraft $w$ of each cloud size: keep the sizes where $w$ beats a pilot's sink rate. Only the two land cases, SCMS and ARM, are used.}
 \figat{8}{24.5}{60}{45}{\figdir/fig7b.png}
 % overlay in Fig. 7b axes: x 0--2500 m on px 120--701, w 0--3 m/s on px 606--20; image 760x735 px drawn 45 mm high
 \fill[black!20,fill opacity=.6](16.06,25.72)rectangle(19.62,61.60);
 \node[rotate=90,font=\fontsize{6}{7}\selectfont,text=black!75]at(17.84,38){thermals 50--300 m};
 \draw[expert,line width=.6pt](15.35,48.45)--(50.92,48.45);
 \draw[pisablue,line width=.6pt](15.35,50.84)--(50.92,50.84);
 \fill[expert,fill opacity=.75](23.88,57.4)rectangle(42.38,58.6);
 \fill[pisablue,fill opacity=.75](21.04,58.9)rectangle(43.80,60.1);
 \figcap{8}{70.3}{60}{Fig. 7b; overlay ours.}
 \head{57}{24.5}{96}{Climbing needs $w$ above the sink rate}
 \smalltxt{57}{30}{96}{$w$: mean updraft of clouds of size $l$. Sink rates are for straight flight: a circling wing sinks faster, so \textbf{the real min.\ sink thresholds are higher}.}
 \txt[font=\fontsize{8.5}{11}\selectfont]{57}{40}{96}{\begin{tabular}{@{}l@{\hspace{3mm}}c@{\hspace{3mm}}c@{\hspace{3mm}}c@{}}
  & min.\ sink & clouds with $w\ge$ sink & thermal diameter\\\hline
  {\color{expert!80!black}\textbf{Paraglider}} & 1.1 m/s & 0.6--1.9 km & 50--300 m\\
  {\color{pisablue}\textbf{Hang glider}} & 0.9 m/s & 0.4--2.0 km & 80--300 m
 \end{tabular}}
 \plainbox{57}{55}{93}{\textbf{Clouds of 50--300 m} are the most numerous, but their mean updraft stays below 0.7 m/s: too weak to keep a paraglider or a hang glider climbing.\\[1mm]\textbf{Clouds with enough updraft} are 0.4--2 km wide: several times, up to ten times, the 50--300 m expected for the thermals pilots use.}
 \bottomline{``Climbable clouds'': 0.4--2 km. Expected thermals: 50--300 m}
 \credit{\S4b, Fig.\ 7b, $w$ read off the plot. Sink rates: typical straight-flight minimum sink. Thermal diameters: rough estimate, not from a measurement.}
\notetext{Neggers et al.\ give no statistics of thermals. Figure 7b gives, for each cloud size $l$, the cloud-mean vertical velocity $w$, which grows with size. A pilot climbs only where the air rises faster than the wing sinks, so the minimum sink rate sets a threshold on $w$: about 1.1 m/s for a paraglider and 0.9 m/s for a hang glider, in straight flight. Circling increases the sink rate, so these thresholds are lower bounds and the true climbable ranges would shift to even larger clouds.\par\smallskip Only the two land cases are used. Over the ocean (BOMEX) the surface sensible heat flux is 8 W\,m$^{-2}$, against 150 W\,m$^{-2}$ in SCMS and up to 140 W\,m$^{-2}$ in ARM (Table 1): without a heated ground, the dry thermals pilots use are weak.\par\smallskip Read off the plot, SCMS clouds exceed 0.9 m/s from about 400 m and 1.1 m/s from about 600 m, up to the largest clouds near 1.8 km. ARM clouds exceed 0.9 m/s from about 700 m and 1.1 m/s from about 850 m, up to about 2 km. Taking both cases, the climbable cloud sizes are 0.6--1.9 km for a paraglider and 0.4--2.0 km for a hang glider. The expected thermal diameters are 50--300 m and 80--300 m. Clouds of that size have $w$ well below either sink rate.\par\smallskip The ranges start above the widest expected thermal and reach 2 km. If cloud sizes were read as thermal sizes, the thermals would come out several times too wide: $l$ is the square root of the projected area, which merges all heights (Fig. 3), and a cloud is fed through its base by a narrower thermal. Caveats: $w$ is the in-cloud mean, which includes weak edges and differs from the updraft below cloud base; the statistics are LES snapshots of two land cases. Source: Section 4b, Figure 7b, Table 1.}
\finishslide\end{frame}
% A: sink rate in a banked turn ---------------------------------------------
\begin{frame}[plain]\frametitle{Appendix: sink rate in a banked turn}\startslide{Appendix: sink rate in a banked turn}
 \lead{Why circling sinks faster than straight flight, at the same angle of attack.}
 % side view, straight glide; path drawn 15 deg below horizontal
 \draw[muted,line width=.4pt](8,24.1)--(55,36.7);
 \draw[-{Latex},black,line width=.8pt](30,30)--(41.6,33.1);
 \draw[-{Latex},teal,line width=.8pt](30,30)--(22.3,27.9);
 \draw[-{Latex},pisablue,line width=.8pt](30,30)--(32.3,21.3);
 \draw[-{Latex},wine,line width=.8pt](30,30)--(30,38);
 \draw[muted,dashed,line width=.4pt](30,30)--(41.6,30);
 \draw[-{Latex},black,line width=.6pt](41.6,30)--(41.6,33.1);
 \fill[black](30,30)circle(1.1);
 \node[anchor=west,text=pisablue,font=\fontsize{7.5}{8.5}\selectfont]at(32.7,22.5){lift $L$};
 \node[anchor=south,text=teal,font=\fontsize{7.5}{8.5}\selectfont]at(20.5,26.4){drag $D$};
 \node[anchor=north,font=\fontsize{7.5}{8.5}\selectfont]at(36.5,32.7){speed $V$};
 \node[anchor=west,font=\fontsize{7.5}{8.5}\selectfont]at(42,31.4){sink $s$};
 \node[anchor=east,text=wine,font=\fontsize{7.5}{8.5}\selectfont]at(29.5,36.6){weight $W$};
 \node[anchor=north west,text=muted,font=\fontsize{6.5}{7.5}\selectfont]at(7,39.2){side view, straight glide};
 % back view, turn to the left; bank 30 deg, lift arrow 13 mm
 \draw[muted,densely dotted](33,46)--(33,72);
 \draw[muted,dashed,line width=.4pt](11,59)--(33,59);
 \fill[muted](11,59)circle(.8);
 \node[anchor=north,text=muted,font=\fontsize{6.5}{7.5}\selectfont,align=center]at(11,60.2){centre\\of turn};
 \draw[black,line width=1.6pt](23.47,64.5)--(42.53,53.5);
 \fill[black](33,59)circle(1.1);
 \draw[muted,dashed,line width=.4pt](26.5,47.74)--(33,47.74);
 \draw[muted,dashed,line width=.4pt](26.5,47.74)--(26.5,59);
 \draw[-{Latex},pisablue,line width=.8pt](33,59)--(26.5,47.74);
 \draw[-{Latex},pisablue!70,line width=.6pt](33,59)--(33,47.74);
 \draw[-{Latex},pisablue!70,line width=.6pt](33,59)--(26.5,59);
 \draw[-{Latex},wine,line width=.8pt](33,59)--(33,70.26);
 \draw[black!70,line width=.4pt](33,53.5) arc[start angle=270,end angle=240,radius=5.5];
 \node[font=\fontsize{7.5}{8.5}\selectfont]at(31.1,51.9){$\varphi$};
 \node[anchor=south east,text=pisablue,font=\fontsize{7.5}{8.5}\selectfont]at(26.5,47.74){lift $L$};
 \node[anchor=west,text=pisablue,font=\fontsize{7}{8}\selectfont]at(33.4,50.5){$L\cos\varphi$};
 \node[anchor=south,text=pisablue,font=\fontsize{6.5}{7.5}\selectfont]at(29.75,59.1){$L\sin\varphi$};
 \node[anchor=west,text=wine,font=\fontsize{7.5}{8.5}\selectfont]at(33.4,68.6){weight $W$};
 \node[anchor=west,text=muted,font=\fontsize{6.5}{7.5}\selectfont]at(42.9,53.5){wing, bank $\varphi$};
 \node[anchor=north west,text=muted,font=\fontsize{6.5}{7.5}\selectfont]at(7,71.8){back view, turn};
 % derivation: step text on the left, its equation on the right
 \txt[font=\fontsize{8.5}{10.3}\selectfont]{60}{22}{57}{{\color{pisablue}\bfseries 1\enspace Energy.} Lift does no work; drag uses $D\,V$ per second, paid for by lost height}
 \txt[font=\fontsize{9.5}{11.5}\selectfont]{122}{22.5}{31}{$W\,s=D\,V$}
 \txt[font=\fontsize{8.5}{10.3}\selectfont]{60}{31}{57}{{\color{pisablue}\bfseries 2\enspace Glide ratio} $E=L/D$ is set by the angle of attack, so $D=L/E$}
 \txt[font=\fontsize{9.5}{11.5}\selectfont]{122}{31.5}{31}{$s=V\,L/(E\,W)$}
 \txt[font=\fontsize{8.5}{10.3}\selectfont]{60}{40}{57}{{\color{pisablue}\bfseries 3\enspace Straight flight:} lift $\approx$ weight, as the path is only $\approx6^\circ$ below horizontal}
 \txt[font=\fontsize{9.5}{11.5}\selectfont]{122}{40}{31}{$L\approx W,\ \ s_0\approx V_0/E$}
 \txt[font=\fontsize{8.5}{10.3}\selectfont]{60}{48.5}{57}{{\color{pisablue}\bfseries 4\enspace Turn:} the vertical part of lift holds the weight, $L\cos\varphi=W$}
 \txt[font=\fontsize{9.5}{11.5}\selectfont]{122}{49}{31}{$L=W/\cos\varphi$}
 \txt[font=\fontsize{8.5}{10.3}\selectfont]{60}{57}{57}{{\color{pisablue}\bfseries 5\enspace Same angle of attack:} lift grows as $V^2$}
 \txt[font=\fontsize{9.5}{11.5}\selectfont]{122}{57}{31}{$V=V_0/\sqrt{\cos\varphi}$}
 \plainbox{60}{62.5}{90}{\fontsize{9.5}{11.5}\selectfont{\color{pisablue}\bfseries Put 3, 4, 5 into 2:}\quad$s=\dfrac{V_0}{\sqrt{\cos\varphi}}\cdot\dfrac{W}{\cos\varphi}\cdot\dfrac{1}{E\,W}=\dfrac{s_0}{(\cos\varphi)^{3/2}}$}
 \bottomline{Circling sinks faster: $\times1.10$ at $20^\circ$ of bank, $\times1.24$ at $30^\circ$, $\times1.68$ at $45^\circ$}
 \credit{Steady flight; small glide angle, so $L\approx W$ (error below 1\% for a glide ratio of 9). The extra drag of turning is neglected, so the true factor is larger.}
\notetext{Derivation of the sink rate in a steady banked turn, used on the previous slide. Step 1: in a steady glide the speed is constant and the lift, perpendicular to the velocity, does no work; the drag dissipates $D V$ per second, and the only source paying for it is height, so $W s=D V$. Step 2: the glide ratio $E=L/D$ depends only on the angle of attack, so $D=L/E$ and $s=V L/(E W)$. Step 3: in straight flight the lift equals the weight, $L=W$, so $s_0=V_0/E$.\par\smallskip Step 4: in a turn banked by $\varphi$ the lift is tilted towards the centre of the turn. Its vertical part must still hold the weight, $L\cos\varphi=W$, so $L=W/\cos\varphi$; its horizontal part, $L\sin\varphi$, bends the path. Step 5: at the same angle of attack the lift grows as the square of the speed, $L=\tfrac12\rho V^2 S C_L$, so a lift $1/\cos\varphi$ times larger needs $V=V_0/\sqrt{\cos\varphi}$, and $E$ is unchanged. Putting 3, 4 and 5 into 2 gives $s=s_0/(\cos\varphi)^{3/2}$: 10\% more at 20$^\circ$ of bank, 24\% at 30$^\circ$, 68\% at 45$^\circ$.\par\smallskip The factor does not depend on the angle of attack: at every angle of attack the turn multiplies the straight-flight sink by $1/(\cos\varphi)^{3/2}$, so the minimum sink in a turn is the straight-flight minimum sink times the same factor. Steps 3 and 4 use $L\approx W$, valid for a small glide angle: with a glide ratio of 9 the path is about 6$^\circ$ below horizontal and the error is below 1\%. Real wings, paragliders in particular, also lose more to turning drag. The straight-flight minimum sink is therefore a lower bound on the updraft needed to climb while circling. The two force balances also give the turn radius, $r=V^2/(g\tan\varphi)$: circling in a narrow thermal needs a steep bank, hence a higher sink rate.}
\finishslide\end{frame}
\end{document}
